\chapter{One and the same boundary: interior flow and gauge tension}
\label{chap:hoja-comun-ns-ym}
\index{Navier--Stokes}\index{Yang--Mills}
\index{boundary operator}\index{Stokes identity}

Hydrodynamics and Yang--Mills are not brought together by externally comparing two
celebrated equations. They follow from two representations of the same boundary
operator. The orientation that the interior reader interprets as circulation
becomes holonomy under the reflected reader; the symmetric memory
that controls the flow becomes a positive transfer kernel.

\section{The common boundary state}

Let \(K_m\) be a finite oriented complex, \(B_m\) its algebra of boundary
data, and \(d_m\) its coboundary. The enriched state is
\[
 s_m=(b_m,U_m,\theta_m,c_m,\Omega_m),
\]
where \(U_m\) transports, \(\theta_m^2=I\) reverses orientation,
\(c_m\) preserves memory, and \(\Omega_m\) is a \(1\)-connection cochain.
The parity projectors are
\[
 E_\pm=\frac12(I\pm\theta_m).
\]
The part \((U_m+U_m^\dagger)/2\) is symmetric; it is a projector only if its
idempotence is also proved.

On the cochain complex, two linear readers are constructed:
\[
 \rho_{{\rm int},m}:C^0(K_m)\longrightarrow
 C^1(K_m)\oplus C^0(K_m),
\qquad
 \rho_{{\rm int},m}(v)=(d_mv,d_m^\ast d_mv),
\]
\[
 \rho_{{\rm ref},m}:C^1(K_m)\longrightarrow
 C^0(K_m)\oplus C^1(K_m),
\qquad
 \rho_{{\rm ref},m}(a)=(d_m^\ast a,d_md_m^\ast a).
\]
The former transports a vertex potential to its edge flux and interior
response; the latter transports an edge cochain to its boundary response and
reflected operator. The codomains differ, but both are constructed with the
same coboundary. Circulation, vorticity, holonomy, and curvature arise
subsequently through Stokes' identity; they are not presupposed within the
symbols \(\rho\).

\section{The minimal operator and its spectrum}

In the oriented cycle of three vertices, let
\[
 d_0=
 \begin{pmatrix}
 -1&1&0\\
 0&-1&1\\
 1&0&-1
 \end{pmatrix}.
\]
The Hodge Laplacians of vertices and edges are
\[
 \mathcal L_{\rm H}^{(0)}=d_0^\ast d_0,
 \qquad
 \mathcal L_{\rm H}^{(1)}=d_0d_0^\ast,
\]
and a direct calculation gives
\[
 \mathcal L_{\rm H}^{(0)}=\mathcal L_{\rm H}^{(1)}=
 \begin{pmatrix}
 2&-1&-1\\
 -1&2&-1\\
 -1&-1&2
 \end{pmatrix}
 =3I-J.
\]
In this residence, the superscripts \((0)\) and \((1)\) denote the
Hodge Laplacians on vertices and edges. They do not denote the
lifting operators \(L_0,L_1\) in the chain \(w_6\to w_{30}\). For
\(\mathcal L_{\rm cyc}^{\rm TPK}:=2I-C-C^{-1}=3I-J\), the realization of the
TPK cyclic operator is typed through
\[
 \mathfrak R_{\rm hyd}(\mathcal L_{\rm cyc}^{\rm TPK})
 =\mathcal L_{\rm H}^{(0)},\qquad
 \mathfrak R_{\rm YM}(\mathcal L_{\rm cyc}^{\rm TPK})
 =\mathcal L_{\rm H}^{(1)}.
\]

\begin{theorem}[Common boundary operator]
\label{thm:operador-comun-frontera-ns-ym}
In the three-cell cycle,
\[
 \operatorname{Spec}(\mathcal L_{\rm H}^{(0)})
 =\operatorname{Spec}(\mathcal L_{\rm H}^{(1)})
 =\{0,3,3\}.
\]
After removing the constant mode ---zero mean in the interior reader and
gauge in the connection reader--- both sectors have exact gap \(3\).
In every finite complex, \(d^\ast d\) and \(dd^\ast\) have the same nonzero
spectrum, with multiplicities.
\end{theorem}

\begin{proof}
The vector \((1,1,1)\) generates the kernel of \(3I-J\). On its orthogonal
complement, \(J=0\), so the operator is \(3I\). For the
general statement, a singular decomposition \(d=U\Sigma V^\ast\) yields
\[
 d^\ast d=V\Sigma^\ast\Sigma V^\ast,
 \qquad
 dd^\ast=U\Sigma\Sigma^\ast U^\ast.
\]
The squares of the non-zero singular values, including their multiplicities, coincide.
\end{proof}

\section{Hydrodynamic reading}

We define the implicit step
\[
 S_{\rm NS}=(I+\mathcal L_{\rm H}^{(0)})^{-1}.
\]
In the minimum cycle,
\[
 S_{\rm NS}=\frac14(I+J).
\]

\begin{theorem}[Global contraction of the finite flow]
\label{thm:contraccion-flujo-finito}
For every mean-zero \(v_0\), the iteration
\(v_{k+1}=S_{\rm NS}v_k\) exists and is unique for every \(k\geq0\), and
\[
 v_k=4^{-k}v_0,
 \qquad
 \|v_k\|\leq4^{-k}\|v_0\|,
 \qquad
 \|v_k\|^2=16^{-k}\|v_0\|^2.
\]
\end{theorem}

\begin{proof}
At zero mean one has \(Jv=0\), hence
\(S_{\rm NS}v=v/4\). The formula is obtained by induction.
\end{proof}

\Needspace{4\baselineskip}
An oriented nonlinearity \(N_m(v,c)\) preserves energy when
\[
 \langle v,N_m(v,c)\rangle=0.
\]
Advection then redistributes the energy, and the positive operator controls
its propagation. The memory \(c_m\) is not absorbed into a chosen viscosity:
it is part of the state that must be transported between scales.

\begingroup
\setlength{\parskip}{0pt}
\titlespacing*{\section}{0pt}{8pt plus 1pt minus .5pt}{2.5pt}
\setlength{\abovedisplayskip}{.25\baselineskip}
\setlength{\belowdisplayskip}{.25\baselineskip}
\setlength{\abovedisplayshortskip}{.12\baselineskip}
\setlength{\belowdisplayshortskip}{.18\baselineskip}
\section{Yang--Mills reading}

A \(1\)-cochain \(a\) is interpreted as a connection potential and has
quadratic energy
\[
 Q_{\rm YM}(a)=\langle a,\mathcal L_{\rm H}^{(1)}a\rangle.
\]

\begin{theorem}[Finite gauge coercivity and decay]
\label{thm:coercividad-gap-gauge-finito}
In the quotient by the constant gauge mode,
\[
 Q_{\rm YM}(a)\geq3\|a\|^2.
\]
The semigroup \(K_t=e^{-t\mathcal L_{\rm H}^{(1)}}\) satisfies
\[
 \|K_ta\|\leq e^{-3t}\|a\|
\]
in the physical sector of the minimal cycle.
\end{theorem}

\begin{proof}
On the complement of the kernel, the
\cref{thm:operador-comun-frontera-ns-ym} gives
\(\mathcal L_{\rm H}^{(1)}=3I\). Both formulas
follow by functional calculus.
\end{proof}

For a non-abelian connection,
\[
 F_\Omega=d\Omega+\Omega\smile\Omega,
\qquad
 W_\gamma(\Omega)=\operatorname{Tr}\operatorname{Hol}_\gamma(\Omega).
\]
Reflection positivity requires a factorization on the positive half,
\[
 \langle\theta f,Kf\rangle=\|Cf\|^2\geq0;
\]
pointwise positivity of the matrix \(K\) does not replace it.

\section[Stokes Identity: Circulation--Vorticity and Holonomy--Curvature]{Discrete Stokes identity and \(2\) realizations: circulation--vorticity and holonomy--curvature}

For a \(1\)-cochain \(\Omega\) and a \(2\)-chain \(K\),
\[
 \sum_{e\subset\partial K}\Omega(e)
 =
 \sum_{f\in K^{(2)}}(d\Omega)(f).
\]
Interior edges cancel by orientation. The hydrodynamic
reader interprets the left-hand side as circulation and the right-hand side
as accumulated vorticity. The gauge reader interprets the former as
linearized holonomy and the latter as curvature.

\begin{theorem}[Common Boundary Sheet Theorem]
\label{thm:hoja-comun-frontera}
The decomposition \(B_m=E_-B_m\oplus E_+B_m\) has two typed
representations:
\[
\boxed{
\begin{array}{ccl}
E_-B_m&\longmapsto&
\text{circulation}\ /\ \text{holonomy},\\[2mm]
E_+B_m&\longmapsto&
\text{dissipative control}\ /\ \text{reflected kernel}.
\end{array}}
\]
Stokes equality and the coincidence of the nonzero spectrum belong to the
common object. The hydrodynamic contraction and the gauge gap are distinct
consequences in their respective codomains.
\end{theorem}

\begin{proof}
The projectors \(E_\pm\) are complementary because
\(\theta_m^2=I\). The Stokes identity provides the intertwining of the
oriented sector; \cref{thm:operador-comun-frontera-ns-ym} provides the
common spectrum. The
\cref{thm:contraccion-flujo-finito,thm:coercividad-gap-gauge-finito}
compute the two realizations.
\end{proof}

The theorem is a finite HMT closure: it neither identifies a PDE with a gauge theory nor confuses the gap \(3\) with a physical mass. It explains by construction why interior regularity and boundary separation are two functions of the same oriented memory. Any subsequent continuous passage must preserve the operator, norm, and spectrum that make this identity exact. \endgroup
